% GENERATED FILE. Edit canonical lessons, then run npm run generate:books.
% canonical-source-hash: fnv1a64:8ceec18370424ca9
% canonical-lessons: ZH-C108-shi, ZH-W108-shi-try, ZH-C108-huan-change, ZH-C108-shuang, ZH-C108-daxiao, ZH-C108-heshi

\chapter{Trying It On}
\label{ch:zh-trying-it-on}
\hlchaptermodality{108}

\begin{chapteropening}
\textbf{By the end of this chapter:} \emph{I can ask to try something on, change it, and talk about pairs, sizes and fit.}

Complete the last lesson of chapter 108: I can ask to try something on, change it, and talk about pairs, sizes and fit.
\end{chapteropening}

\section[\texorpdfstring{\zh{试} (shì)}{试 (shì)}]{\zh{试} (shì) — to try; to try on}

\label{lesson:ZH-C108-shi}



\begin{quote}
Before the new one: say the Mandarin for to stop, then the Mandarin for to leave a place.
\end{quote}

\subsection*{You'll want to know: \zh{试}}
\textbf{\zh{试}} — \emph{shì} — \textquotedblleft{}to try; to try on\textquotedblright{}.

\begin{quote}
\textbf{\zh{这个能试吗}} \emph{zhège néng shì ma} — Can I try this one on?
\end{quote}

The new character, in pieces: \zh{试} — the speech piece of \zh{说} on the left; on the right six strokes: a horizontal across the top, a small three-stroke piece low under it, and a long stroke that slants down and hooks, with a dot at its top right. The next lesson puts it in your hand.

\begin{cousinweb}
\textbf{\zh{试}} (\emph{shì}) is inherited from Middle Chinese \emph{syiH} (Guangyun, in Baxter's transcription), glossed 'to use'. Trying a thing by putting it to use is the usual explanation, but the step is not given there.
\end{cousinweb}

\begin{grammarlens}[title={What you've built}]
One more word for asking for what I want.
\end{grammarlens}

\subsection*{Your turn}
\begin{itemize}
\raggedright
  \item \emph{Say it:} \emph{shì}
  \item \emph{Say it:} \emph{shì}, once more
  \item \emph{Say it:} \emph{zhège néng shì ma} — Can I try this one on?
  \item \emph{Recall:} say the Mandarin for to stop, then the Mandarin for to leave a place, then say \emph{shì} again
\end{itemize}

\begin{tcolorbox}[breakable,colback=teal!4,colframe=teal!35!black,title={Before you move on}]
What does \zh{试} mean? (\textquotedblleft{}To try; to try on\textquotedblright{}.) Say it once more.
\end{tcolorbox}
\section[\texorpdfstring{\zh{试} (shì)}{试 (shì)}]{\zh{试} — the speech piece, then a slanting hook}

\label{lesson:ZH-W108-shi-try}



\begin{quote}
The word first, then the character inside it.

\begin{itemize}
\raggedright
  \item \emph{Recall:} say \emph{shì} — \textbf{R1}, one lesson back, before it decays
\end{itemize}

\begin{quote}
\zh{说} \zh{请}
\end{quote}

\emph{Read it:} each one aloud, with nothing said first — \emph{shuō}, \emph{qǐng}
\end{quote}

\subsection*{Script — the speech piece, then a slanting hook}
\hlblockfigure{\detokenize{figures/ZH-W108-shi-try-filmstrip.pdf}}{How \zh{试} is written, stroke by stroke}

\begin{quote}
\zh{试}
\end{quote}

\textbf{Two pieces, left before right.} Count the pieces, not the strokes.

\textbf{The speech piece, two strokes}, as in \zh{说}: a dot, then a stroke across that turns down and rises.

\textbf{The right piece, six strokes}: a horizontal across the top; then, low on the left, a short horizontal, a vertical and a stroke that rises to the right; then one long stroke that slants down to the right and hooks up; and last a dot at the top right.

2 + 6 = eight.

\subsection*{Your turn}
\begin{itemize}
\raggedright
  \item \emph{Write it:} \zh{试}, all eight, with the model in front of you
  \item \emph{Cover:} write it again from memory, and count the pieces, not the strokes
\end{itemize}

\begin{tcolorbox}[breakable,colback=teal!4,colframe=teal!35!black,title={Before you move on}]
How many pieces in \zh{试}? (\textbf{Two} — the speech piece, then the right piece.) Which characters you know start with the same piece? (\textbf{\zh{说}}, \textbf{\zh{请}}, \textbf{\zh{课}}.) Which stroke comes last? (\textbf{The dot at the top right}.)

Source: \href{https://raw.githubusercontent.com/chanind/hanzi-writer-data/68d10a4b21150cae5e1ebbd223eed289cf32d90c/data/\%E8\%AF\%95.json}{Hanzi Writer Data \zh{试}, pinned stroke paths and medians}.
\end{tcolorbox}
\section[huàn]{huàn — to change; to exchange}

\label{lesson:ZH-C108-huan-change}



\begin{quote}
Before the new one: say the Mandarin for to leave a place, then the Mandarin for to try on.
\end{quote}

\subsection*{You'll want to know: huàn}
\emph{huàn} — \textquotedblleft{}to change; to exchange\textquotedblright{}.

\begin{quote}
\emph{wǒ xiǎng huàn yí ge} — I'd like to change it for another one.
\end{quote}

Its characters come later; for now it is a word for the ear.

\begin{cousinweb}
\emph{huàn} is inherited from Middle Chinese \emph{hwanH} (Guangyun, in Baxter's transcription), glossed 'to exchange'.
\end{cousinweb}

\begin{grammarlens}[title={What you've built}]
One more word for asking for what I want.
\end{grammarlens}

\subsection*{Your turn}
\begin{itemize}
\raggedright
  \item \emph{Say it:} \emph{huàn}
  \item \emph{Say it:} \emph{huàn}, once more
  \item \emph{Say it:} \emph{wǒ xiǎng huàn yí ge} — I'd like to change it for another one.
  \item \emph{Recall:} say the Mandarin for to leave a place, then the Mandarin for to try on, then say \emph{huàn} again
\end{itemize}

\begin{tcolorbox}[breakable,colback=teal!4,colframe=teal!35!black,title={Before you move on}]
What does huàn mean? (\textquotedblleft{}To change; to exchange\textquotedblright{}.) Say it once more.
\end{tcolorbox}
\section[shuāng]{shuāng — a pair of (a measure word)}

\label{lesson:ZH-C108-shuang}



\begin{quote}
Before the new one: say the Mandarin for to try on, then the Mandarin for to exchange.
\end{quote}

\subsection*{You'll want to know: shuāng}
\emph{shuāng} — \textquotedblleft{}a pair of (a measure word)\textquotedblright{}.

\begin{quote}
\emph{wǒ xiǎng mǎi yì shuāng xié} — I'd like to buy a pair of shoes.
\end{quote}

Its characters come later; for now it is a word for the ear.

\begin{cousinweb}
\emph{shuāng} is inherited from Middle Chinese \emph{sraewng} (Guangyun, in Baxter's transcription), glossed 'a pair, two of a kind'.
\end{cousinweb}

\begin{grammarlens}[title={What you've built}]
One more word for asking for what I want.
\end{grammarlens}

\subsection*{Your turn}
\begin{itemize}
\raggedright
  \item \emph{Say it:} \emph{shuāng}
  \item \emph{Say it:} \emph{shuāng}, once more
  \item \emph{Say it:} \emph{wǒ xiǎng mǎi yì shuāng xié} — I'd like to buy a pair of shoes.
  \item \emph{Recall:} say the Mandarin for to try on, then the Mandarin for to exchange, then say \emph{shuāng} again
\end{itemize}

\begin{tcolorbox}[breakable,colback=teal!4,colframe=teal!35!black,title={Before you move on}]
What does shuāng mean? (\textquotedblleft{}A pair of (a measure word)\textquotedblright{}.) Say it once more.
\end{tcolorbox}
\section[\texorpdfstring{\zh{大小} (dàxiǎo)}{大小 (dàxiǎo)}]{\zh{大小} (dàxiǎo) — size}

\label{lesson:ZH-C108-daxiao}



\begin{quote}
Before the new one: say the Mandarin for to exchange, then the Mandarin for a pair of.
\end{quote}

\subsection*{You'll want to know: \zh{大小}}
\textbf{\zh{大小}} — \emph{dàxiǎo} — \textquotedblleft{}size\textquotedblright{}.

\begin{quote}
\textbf{\zh{我试了大小很好}} \emph{wǒ shì le dàxiǎo hěn hǎo} — I've tried it on, and the size is fine.
\end{quote}

\begin{cousinweb}
\textbf{\zh{大小}} (\emph{dàxiǎo}) is built from \zh{大} 'big' and \zh{小} 'small', both inherited from Middle Chinese, \emph{daH} and \emph{sjewX} (Guangyun, in Baxter's transcription): how big or small a thing is.
\end{cousinweb}

\begin{grammarlens}[title={What you've built}]
One more word for asking for what I want.
\end{grammarlens}

\subsection*{Your turn}
\begin{itemize}
\raggedright
  \item \emph{Say it:} \emph{dàxiǎo}
  \item \emph{Say it:} \emph{dàxiǎo}, once more
  \item \emph{Say it:} \emph{wǒ shì le dàxiǎo hěn hǎo} — I've tried it on, and the size is fine.
  \item \emph{Recall:} say the Mandarin for to exchange, then the Mandarin for a pair of, then say \emph{dàxiǎo} again
\end{itemize}

\begin{tcolorbox}[breakable,colback=teal!4,colframe=teal!35!black,title={Before you move on}]
What does \zh{大小} mean? (\textquotedblleft{}Size\textquotedblright{}.) Say it once more.
\end{tcolorbox}
\section[héshì]{héshì — right; suitable; a good fit}

\label{lesson:ZH-C108-heshi}



\begin{quote}
Before the new one: say the Mandarin for a pair of, then the Mandarin for size.
\end{quote}

\subsection*{You'll want to know: héshì}
\emph{héshì} — \textquotedblleft{}right; suitable; a good fit\textquotedblright{}.

\begin{quote}
\emph{zhè shuāng xié hěn héshì} — These shoes fit well.
\end{quote}

Its characters come later; for now it is a word for the ear.

\begin{cousinweb}
\emph{héshì} is built from hé and shì, both inherited from Middle Chinese, \emph{hop} and \emph{syek} (Guangyun, in Baxter's transcription); the dictionary glosses the first 'to agree, to match' and the second 'good', among other senses.
\end{cousinweb}

\begin{grammarlens}[title={What you've built}]
One more word for asking for what I want.
\end{grammarlens}

\subsection*{Your turn}
\begin{itemize}
\raggedright
  \item \emph{Say it:} \emph{héshì}
  \item \emph{Say it:} \emph{héshì}, once more
  \item \emph{Say it:} \emph{zhè shuāng xié hěn héshì} — These shoes fit well.
  \item \emph{Recall:} say the Mandarin for a pair of, then the Mandarin for size, then say \emph{héshì} again
\end{itemize}

\begin{tcolorbox}[breakable,colback=teal!4,colframe=teal!35!black,title={Before you move on}]
What does héshì mean? (\textquotedblleft{}Right; suitable; a good fit\textquotedblright{}.) Say it once more.
\end{tcolorbox}
